\documentclass[a4paper]{scrreprt} % Рекомендований клас для навчальних видань

% lecture_onu.sty
% Стильовий файл, адаптований до вимог ОНУ (14 кегель, 1.5 інтервал, поля 2 см)
\ProvidesPackage{lecture_onu}

% !TeX program = xelatex
% !TeX options = -shell-escape
\usepackage{ragged2e}
\usepackage{cmap}
\usepackage[ukrainian]{babel}

% 1. МОВА: Встановлено українську мову
\usepackage{algorithm}
\usepackage{algorithmic}


%tabularx, бо він сам розподіляє простір і не дає таблиці «вилазити» за поля.
\usepackage{tabularx}

% 3, 4. ПОЛЯ ТА РОЗМІР: Встановлено поля 2 см та розмір шрифту 14 пт (для KOMA-script)
\KOMAoptions{fontsize=14pt}
\usepackage[left=2cm, right=2cm, top=2cm, bottom=2cm]{geometry}


\usepackage{iftex}
\ifPDFTeX
\usepackage[
activate={true,nocompatibility},
final,
tracking=true,
factor=1100,
stretch=10,
shrink=10,
letterspace=700
]{microtype}
\else
\usepackage[stretch=10,shrink=10,tracking=true]{microtype}
\fi


% 2. ШРИФТ: Видалено kpfonts. Встановлюємо Times New Roman.
\usepackage{fontspec}
% ПРИМІТКА: Переконайтеся, що "Times New Roman" встановлений у вашій системі. 
% Для Linux-систем може знадобитися Tinos або Liberation Serif.
\setmainfont{Times New Roman}
\setsansfont{Times New Roman}
\setmonofont{Courier New} % Змінено на більш стандартний моноширинний

%автоматично коригує інтервали між буквами та словами
\usepackage{microtype}


% Видалено mathastext та пов'язані пакети kpfonts для дотримання формату
\usepackage{amsmath,amssymb,amsfonts}
\usepackage{icomma}
\usepackage{mathtools}
\usepackage{nicefrac}
\usepackage{amsthm}


% 5. ІНТЕРВАЛ ТА ВІДСТУП: Встановлено 1.5 інтервал. Видалено parskip.
\usepackage{setspace}
\onehalfspacing 
\setlength{\parskip}{0pt} % Видалено додатковий інтервал між абзацами
\setlength{\parindent}{1cm} % Залишено стандартний абзацний відступ

\usepackage{enumitem}
\setlist{noitemsep} % Залишено, для охайних списків
\renewcommand{\labelitemi}{\faCheckSquareO}

\usepackage[bottom]{footmisc}
\setlength{\footnotesep}{12pt}%

\usepackage[
automark,
headsepline,
footsepline,
autooneside
]{scrlayer-scrpage}
\renewcommand*{\chaptermarkformat}{}
\renewcommand*{\sectionmarkformat}{}
\automark[]{chapter}
\ihead[]{\leftmark}
\chead[]{}
\ohead[]{}




% 7. КОПІРАЙТ: Змінено на українську мову
\ifoot[\textcopyright~{Платонова Є. В., Страхов Є. М.}]{\textcopyright~{ Платонова Є. В., Страхов Є. М.}}
\cfoot[]{}
\ofoot[\pagemark]{\pagemark}
\setkomafont{pageheadfoot}{\footnotesize}
\setkomafont{pagehead}{\small\bfseries\itshape}
\setkomafont{pagenumber}{\normalsize}
\pagestyle{scrheadings}

% СТИЛІЗАЦІЯ ЗАГОЛОВКІВ (Залишено, але пам'ятайте, що вона не є стандартною для ДСТУ)
\usepackage{titlesec}
\titleformat{\chapter}[display]
{\sffamily\bfseries}
{\begin{tikzpicture}[remember picture,overlay]
		\node[below left, shade, top color=black, bottom color=pyblue, rectangle, inner sep=2pt, xshift=-2.0cm, minimum width=1.75cm, minimum height=5cm, text width=1.75cm, align=center] at (current page.north east) {\color{pyyellow}\fontsize{50}{60}\selectfont\mbox{}\\[2cm]\thechapter};
\end{tikzpicture}}
{-1.425cm}
{\Huge}

\titlespacing*{\chapter}
{0pt}{0pt}{30pt}

\usepackage{fancyvrb}
\usepackage[dvipsnames,svgnames,x11names]{xcolor}
\usepackage{minted}
\usemintedstyle{vs}
\usepackage[many]{tcolorbox}
\tcbuselibrary{skins}
\tcbuselibrary{minted}
\tcbuselibrary{theorems}

\titleformat{\subsection}[block] % block = окремий рядок
{\bfseries\normalsize}         % жирний + стандартний розмір
{}                             % номер (порожньо для subsection*)
{0pt}                          % відступ
{}   



% Оформлення блоків (прикладів, задач) залишено.
\newtcolorbox[auto counter, number within=chapter]{pr}[1]{
	enhanced, 
	arc=0mm, outer arc=0mm,
	frame hidden, 
	borderline north={1pt}{-1pt}{MidnightBlue}, borderline south={2pt}{-2pt}{MidnightBlue},
	boxrule=0pt, titlerule=0pt,
	left=0.2cm, right=0.2cm,
	toptitle=0.1cm, bottomtitle=0.1cm, top=0.2cm, colbacktitle=MidnightBlue,
	colback=GhostWhite, coltitle=white, before skip=12pt plus 2pt minus 2pt, after skip=10pt plus 2pt minus 2pt, fonttitle=\bfseries, title={Задача~\thetcbcounter}, #1
} % умови задач

\newtcblisting{py}[1]{
	enhanced, arc=0pt, outer arc=0pt, 
	listing engine=minted, minted style=vs, minted language=python, minted options={fontsize=\small, numbersep=3mm, baselinestretch=1.1}, fonttitle=\bfseries, title={#1}, 
	listing only,
	colback=MintCream, colbacktitle=pyblue,
	frame hidden, 
	borderline north={1pt}{-1pt}{pyblue}, borderline south={2pt}{-2pt}{pyblue}, 
	boxrule=0pt, titlerule=0pt,
	left=0.2cm, right=0.2cm,  
	toptitle=0.1cm, bottomtitle=0.1cm, top=0.2cm, before skip=10pt plus 2pt minus 2pt, after skip=10pt plus 2pt minus 2pt, after={\par\noindent}
} % код

\newtcblisting{pydef}{
	enhanced, arc=0pt, outer arc=0pt, 
	listing engine=minted, minted style=vs, minted language=python, minted options={fontsize=\small, numbersep=3mm}, 
	colback=orange!10!white,
	listing only,
	frame hidden, 
	borderline north={1pt}{-1pt}{orange}, borderline south={2pt}{-2pt}{orange}, 
	boxrule=0pt, titlerule=0pt,
	left=0.2cm, right=0.2cm, 
	toptitle=0.1cm, bottomtitle=0.1cm, top=0.2cm, 
	before skip=10pt plus 2pt minus 2pt, after skip=10pt plus 2pt minus 2pt, after={\par\noindent}
} % шпаргалки с кодом

\newtcolorbox{defn}[1]{
	enhanced, arc=0pt, outer arc=0pt, 
	colback=orange!10!white,
	fonttitle=\bfseries, title={#1}, before skip=10pt plus 2pt minus 2pt, after skip=10pt plus 2pt minus 2pt,
	colbacktitle=orange,  
	frame hidden, 
	borderline north={1pt}{-1pt}{orange}, borderline south={2pt}{-2pt}{orange}, 
	boxrule=0pt, titlerule=0pt,
	left=0.2cm, right=0.2cm,  toptitle=0.1cm, bottomtitle=0.1cm, top=0.2cm
} % головне

\newtcolorbox{pyr}{
	enhanced, arc=0pt, outer arc=0pt, 
	colback=white, 
	frame hidden, 
	before skip=10pt plus 2pt minus 2pt, after skip=10pt plus 2pt minus 2pt, 
	boxrule=0pt, 
	left=-0.1cm, right=-0.1cm, titlerule=0em, toptitle=0cm, bottomtitle=0cm, top=0.2cm
} % код + результаты

\newtcolorbox{fl}[1]{
	arc=0pt, outer arc=0pt, colback=Gainsboro!30, colframe=darkgray, left=2mm, enhanced, coltitle=white, colbacktitle=black, 
	fonttitle=\ttfamily\bfseries, title={#1}, 
	frame hidden, 
	borderline north={1pt}{-1pt}{black}, borderline south={2pt}{-2pt}{black},
	boxrule=0pt, titlerule=0pt,
	left=0.2cm, right=0.2cm,  toptitle=0.1cm, bottomtitle=0.1cm, top=0.2cm, fontupper=\small\ttfamily
} % файли

\newtcolorbox{res}[1]{
	enhanced, arc=0pt, outer arc=0pt, 
	fonttitle=\bfseries, title={#1}, fontupper=\small, colback=MintCream, colframe=MediumBlue, listing only, 
	boxrule=0pt, titlerule=0pt,
	frame hidden, 
	borderline north={1pt}{-1pt}{MediumBlue}, borderline south={2pt}{-2pt}{MediumBlue},
	left=0.2cm, right=0.2cm,  toptitle=0.1cm, bottomtitle=0.1cm, top=0.2cm, 
	before skip=10pt plus 2pt minus 2pt, after skip=10pt plus 2pt minus 2pt
} % результати

% ... решта пакетів та макросів залишаються, якщо вони не суперечать вимогам
\usepackage{graphicx}
\usepackage{tikz}
\usepackage{pgfplots}




\usepackage{hyperref}
\hypersetup{
	pdftitle={Конспект лекцій}, % Змінено заголовок
	pdfauthor={Євген Страхов, Євгенія Платонова}, % Змінено автора
	unicode=true,
	colorlinks=true,
	bookmarks=true,
	linkcolor=MediumBlue,
	citecolor=FireBrick,
	urlcolor=DodgerBlue
}

\usepackage{menukeys}
\usepackage[titles]{tocloft}
\renewcommand{\cftchappagefont}{\rmfamily}
\renewcommand{\cftpartpagefont}{\rmfamily}
\setcounter{secnumdepth}{2}

\newcommand{\mi}[1]{\mintinline{python}{#1}}
\usepackage{python}
\usetikzlibrary{fadings}
\newcommand\fadingtext[3][]{%
	\begin{tikzfadingfrompicture}[name=fading letter]
		\node[text=transparent!0,inner xsep=0pt,outer xsep=0pt,#1] {#3};
	\end{tikzfadingfrompicture}%
	\begin{tikzpicture}[baseline=(textnode.base)]
		\node[inner sep=0pt,outer sep=0pt,#1](textnode){\phantom{#3}}; 
		\shade[path fading=fading letter,#2,fit fading=false]
		(textnode.south west) rectangle (textnode.north east);% 
	\end{tikzpicture}% 
}
\definecolor{pyblue}{HTML}{3C7DA7}
\definecolor{pyyellow}{HTML}{FFDB6D}
\newcommand{\Python}{{\fontspec[Scale=1.17]{Inconsolata}\color{pyblue}python}}
\usepackage{fontawesome}
\newcommand{\myhref}[2]{\href{#1}{\sffamily #2~\faExternalLink}}
\usepackage{booktabs}
\usepackage{makecell}
\usepackage{multirow}
\usepackage{multicol}
\usepackage{longtable}
\RequirePackage[square,numbers,sort&compress]{natbib}
\setlength{\bibsep}{6pt}

% 9. ЗМІСТ: Встановлено українську назву
\addto\captionsukrainian{%
	\renewcommand{\contentsname}{Зміст}%
}

% 8. ТИТУЛЬНИЙ АРКУШ: Змінено текст на український
\newcommand{\mytitlepage}[2]{
	% \ThisTileWallPaper... (Зображення фону, можливо, доведеться видалити для формальності)
	\thispagestyle{empty}
	\vspace*{2cm}
	\begin{tcolorbox}[
		beamer,
		interior style={left color=LemonChiffon,right color=yellow},
		boxrule=0pt,
		colframe=white,
		arc=10pt,
		outer arc=10pt,
		remember as=title,
		leftrule=0.25cm,
		rightrule=0pt,
		fontupper=\color{white}\sffamily\bfseries\huge,
		enlarge left by=-1.2in-\hoffset-\oddsidemargin, 
		width=0.7\paperwidth, 
		left=1cm,
		top=0.3cm, 
		bottom=0.3cm, 
		]
		{\large\bfseries\color{MidnightBlue}\fontspec[Ligatures=TeX]{Exo 2 Semi Bold Expanded}
			Основи програмування. Курс лекцій}
	\end{tcolorbox}
	\vspace{8cm}
	\begin{tcolorbox}[
		beamer,
		interior style={left color=MidnightBlue,right color=DodgerBlue},
		boxrule=3pt,
		colframe=white,
		arc=10pt,
		remember as=title,
		leftrule=3pt,
		rightrule=0pt,
		fontupper=\color{white}\sffamily\bfseries\huge,
		enlarge left by=-\hoffset-\oddsidemargin, 
		enlarge right by=-\paperwidth+1in+\hoffset+\oddsidemargin+\textwidth,
		width=0.9\paperwidth, 
		left=1cm, 
		right=\paperwidth-1in-\hoffset-\oddsidemargin-\textwidth,
		top=0.75cm, 
		bottom=0.75cm, 
		]
		{\noindent\hfill\huge\bfseries\fontspec{Exo 2 Semi Bold Expanded}
			Частина~#1\\[0.1cm]
			\mbox{}\hfill\MakeUppercase{#2}}
	\end{tcolorbox}
	\vfill
	\begin{tcolorbox}[
		beamer,
		interior style={left color=LemonChiffon,right color=yellow},
		boxrule=0pt,
		colframe=white,
		arc=10pt,
		outer arc=10pt,
		remember as=title,
		leftrule=0.25cm,
		rightrule=0pt,
		fontupper=\color{white}\sffamily\bfseries\huge,
		enlarge left by=-1.2in-\hoffset-\oddsidemargin, 
		width=0.7\paperwidth, 
		left=1cm,
		top=0.3cm, 
		bottom=0.3cm, 
		]
		{\normalsize\bfseries\color{MidnightBlue}\fontspec{Exo 2 Semi Bold Expanded}
			Євген Страхов\\[-5\jot]
			Одеський національний університет, 2026} % Змінено рік
	\end{tcolorbox}
}



\begin{document}
	
	% === ТИТУЛЬНА СТОРІНКА ===
	% Необхідно заповнити фактичними даними:
	% \mytitlepage{Частина №}{Підзаголовок}
	%\mytitlepage{1}{Вступ до Програмування} 
	
\titleformat{\section}{\normalfont\Large\bfseries}{\thesection}{1em}{}
	
	
	
	
	% === ОСНОВНИЙ ТЕКСТ ЛЕКЦІЇ ===
	% Вміст із файлу l1.tex:
		%----------------------------------------------------------------------------------------
%	TITLE PAGE
%----------------------------------------------------------------------------------------
\KOMAoptions{fontsize=12pt} 
\begin{titlepage} % Suppresses displaying the page number on the title page and the subsequent page counts as page 1
	\newcommand{\HRule}{\rule{\linewidth}{0.5mm}} % Defines a new command for horizontal lines, change thickness here
	
	\centering % Centre everything on the page
	
	\Large МІНІСТЕРСТВО ОСВІТИ І НАУКИ УКРАЇНИ\\  % Main heading such as the name of your university/college
	\medskip
	\large ОДЕСЬКИЙ НАЦІОНАЛЬНИЙ УНІВЕРСИТЕТ ІМЕНІ І.І. МЕЧНИКОВА% Main heading such as the name of your university/college
	
	\large ФАКУЛЬТЕТ МАТЕМАТИКИ, ФІЗИКИ ТА ІНФОРМАЦІЙНИХ ТЕХНОЛОГІЙ % Major heading such as course name
	
	%\textbf{\Large Кафедра Оптимального Керування та Економічної Кібернетики}\\[1.5cm] % Minor heading such as course title
	
	\vspace*{5\baselineskip} % Whitespace above the title
	
	
	{\Large 	\textbf{Євгенія ПЛАТОНОВА, Євген СТРАХОВ}\\} % Editor list
	
	\vspace{1\baselineskip} % Whitespace below the editor list
	\rule{\textwidth}{1.6pt}\vspace*{-\baselineskip}\vspace*{2pt} % Thick horizontal rule
	\rule{\textwidth}{0.4pt} % Thin horizontal rule
	
	\vspace{1\baselineskip} % Whitespace above the title
	
	{\LARGE \textbf{Програмування}} % Title
	
	\vspace{0.5\baselineskip} % Whitespace below the title
	
	\rule{\textwidth}{0.4pt}\vspace*{-\baselineskip}\vspace{3.2pt} % Thin horizontal rule
	\rule{\textwidth}{1.6pt} % Thick horizontal rule
	
	\vspace{1\baselineskip} % Whitespace after the title block
	
	%------------------------------------------------
	%	Subtitle
	%------------------------------------------------
	
	\LARGE \textbf{Основи програмування на мові Python.\\ Структурне програмування\\} % Subtitle or further description
	%{(адаптація learnopengl.com)\\}
	
	\vspace*{0.5\baselineskip} % Whitespace under the subtitle
	
	\Large \textbf{Конспект лекцій}
	\vspace*{1.5\baselineskip}
	%\vfill\v% Position the date 3/4 down the remaining page
	
	%		Рекомендовано Методичною радою ОНУ імені І.І. Мечникова\\
	%		як навчальний посібник для здобувачів ступеня бакалавра\\ 
	%		спеціальності 113 «Прикладна математика», 122 «Комп’ютерні науки»,\\
	%		123 «Комп’ютерна інженерія», 126 «Інформаційні системи та технології»\\
	
	
	\vspace*{0.25\baselineskip}
	
	%	\texttt{Електронне мережне навчальне видання}
	\texttt{Електронне навчальне видання}
	
	%{\large\today} % Date, change the \today to a set date if you want to be precise
	
	%------------------------------------------------
	%	Logo
	%------------------------------------------------
	
	%------------------------------------------------
	%	Publisher
	%------------------------------------------------
	\vspace*{4\baselineskip}
	{ \large 
		ОДЕСА\\ \vspace{-1.5mm}
		ОНУ\\ \vspace{-2.5mm}
		2026 }%Publication year
	%\vfill\vfill
	%\includegraphics[width=0.2\textwidth]{placeholder.jpg}\\[1cm] % Include a department/university logo - this will require the graphicx package
	
	%----------------------------------------------------------------------------------------
	
	
	\pagebreak
	
	\large
	\raggedright \Large {\textbf{УДК 004.4:004.421 }\\
		\hspace*{1 cm} XXX}
	
	\bigskip
	\justifying
	\large
	\hspace{1cm}\textbf{Автори:}\\
   	\textbf{\textit{Є. В. Платонова}} старший викладач кафедри оптимального керування і економічної кібернетики;\\ 
	\textbf{\textit{Є. М. Страхов}} канд. фіз.-мат. наук, доцент, доцент кафедри оптимального керування і економічної кібернетики.\\ 
	\medskip
	\hspace{1cm}\textbf{Рецензенти:}\\
	\textbf{\textit{В. О. Сперанський }}  канд. техн. наук, доцент, доцент кафедри комп’ютеризовані системи та програмні технології НУ "Одеська Політехніка";\\
	\textbf{\textit{З. Ю. Журавльова}} канд. фіз.-мат. наук, доцентом кафедри методів математичної фізики.
	
	%		\begin{tabbing}
		%			\textbf{Рецензенти: }  \= \hspace{ 1cm} \= \kill
		%			\textbf{Рецензенти:} \> \textit{\textbf{Полякова М.В.}}, докт. техн. наук, проф., професор кафедри\\ \> прикладної математики та інформаційних технологій\\ \> НУ "Одеська Політехніка".\\ 
		%			 \> \textit{\textbf{Батюк А.Є.}}, канд.техн.наук, доцент кафедри АСУ\\ \> Інституту комп'ютерних наук та інформаційних технологій\\ \> НУ "Львівська політехніка". \\
		%			 \> \textit{\textbf{Батюк А.Є.}}, канд.техн.наук, доцент кафедри \\ \> Одеського національного університету імені І. І. Мечникова.\\ \> НУ "Львівська політехніка". 
		%		\end{tabbing}
	%		\begin{tabbing}
		%				Відповідальний \= \hspace{3cm}\= \kill
		%				Відповідальний \> \\
		%				редактор: \> \textit{Тарасенко О.В.}, кандидат техн. наук, доцент
		%		\end{tabbing}
	
	
	
	\bigskip
	\medskip
	
	\centering {\normalsize \textit{Рекомендовано до видання Методичною радою ОНУ імені І.І. Мечникова\\
			(протокол № \underline{\hspace{1cm}} від \underline{\hspace{1cm}}.\underline{\hspace{1cm}}.2026 р.) } } \\
	%		за поданням Вченої ради факультету математики, фізики\\ та інформаційних технологій\\
	%		(протокол № \underline{\hspace{1cm}} від \underline{\hspace{1cm}}.\underline{\hspace{1cm}}.2024 р.)}

\bigskip
\medskip

\hspace{1.25cm} \raggedright \textbf{Платонова Є.В., Страхов Є.М.}\\ 
\textbf{XXX} \hspace{0.25cm}\justifying \indent Основи програмування на мові Python. Структурне програмування [Електронний ресурс]: конспект лекцій для здобув. першого (бакалавр) рівня вищої освіти спец. E7 "Математика" / Є.В. Платонова, Є.М. Страхов --Електронні текстові дані (1 файл: 364 Кб). Одеса: Одес. нац. ун-т ім. І. І. Мечникова, 2026. -- 55 с.
ISBN   XXX-XXX-XXX-XXX-X 

\bigskip
%\parbox{\textwidth}
\justifying 
{\textit {\small Конспект лекцій є введенням у базові принципи програмування мовою Python та систематизацією основних алгоритмічних конструкцій. У матеріалах конспекту розглядаються три базові конструкції управління виконанням програми, надаються приклади реалізації алгоритмів у Python. 
	Конспект лекцій розроблений для студентів першого курсу спеціальності E7 «Математика» та може бути використаний для самостійного навчання. Матеріал спрямований на формування 
	алгоритмічного мислення та набуття практичних навичок програмування мовою Python.}

\raggedleft	\textbf{УДК 004.4:004.421 }

\small
\begin{tabbing}
	ISBN   XXX-XXX-XXX-XXX-X \= \hspace{4cm} \= ©  \= Платонова Є.В., \= Страхов Є.М., 2026 \\
	\>   \hspace{3cm} ©  Одеський національний університет\\  \> \> \> імені І. І. Мечникова, 2026
\end{tabbing}
%	ISBN   XXX-XXX-XXX-XXX-X                                                     ©  Платонова Є. В., 2026
%																										©  Одеський національний університет	імені І. І. Мечникова, 2025
\end{titlepage}

\linespread{1.25}

\KOMAoptions{fontsize=14pt} 


	% === ЗМІСТ ===
	\tableofcontents
	
	\section*{Вступ}

Дисципліна «Програмування» є базовою складовою підготовки студентів спеціальності 
Е7 «Математика» та формує фундаментальні знання і навички, необхідні для подальшого 
застосування сучасних інформаційних технологій у наукових дослідженнях та професійній діяльності. 
Вона забезпечує розвиток алгоритмічного мислення, вміння будувати та аналізувати алгоритми, 
а також практичні навички роботи з мовами програмування.

Особливе місце у курсі займає мова Python, яка завдяки простоті синтаксису, універсальності 
та широкому спектру бібліотек стала одним із ключових інструментів у сфері аналізу даних, 
машинного навчання та Data Science. Для студентів-математиків це відкриває можливості 
поєднання теоретичних знань з практичними застосуваннями у статистиці, моделюванні, 
обробці великих масивів даних та побудові прогнозних моделей.

Основними завданнями курсу є:
\begin{itemize}
	\item ознайомлення студентів з базовими алгоритмічними конструкціями структурного програмування;
	\item формування навичок написання програм мовою Python;
	\item розвиток умінь застосовувати програмування для розв’язання математичних задач;
	\item підготовка до використання Python у подальших дисциплінах, пов’язаних з аналізом даних та Data Science.
\end{itemize}
Таким чином, курс «Програмування» є важливим етапом у професійній підготовці студентів 
спеціальності Е7 «Математика», оскільки створює основу для подальшого опанування сучасних 
методів аналізу даних, статистичних моделей та технологій штучного інтелекту.
	\chapter{Програмування}

\section{Цілі курсу}
\begin{itemize}
	\item навчитися складати алгоритми
	\item навчиться писати фрагменти коду (текст програми)
	\item навчиться читати чужий код і використовувати його.
\end{itemize}

\section{Що таке програмування?}


Слово «програма» походить від грецьких слів «pro» (перед) і «graphein» (писати). Іншими словами, це був акт написання чогось наперед або достроково.

\begin{defn}{Програмування — це}
	процес проектування, написання, тестування, вадження і підтримки комп’ютерних програм. У вужчому значенні програмування розглядається як кодування — реалізація у вигляді програми одного чи кількох взаємопов’язаних алгоритмів.
\end{defn}

У широкому сенсі процес програмування охоплює і створення, тобто розробку, алгоритмів, і аналіз потреб майбутніх користувачів програмного забезпечення. Технологія програмування постійно розвивається, що призводить до появи нових мов програмування, інструментів розробки і сфер застосування. Тому важливою ознакою програмування є регулярне опанування нових технологій.

\begin{defn}{Алгоритм — це}
	скінченна послідовність вказівок на виконання дій, спрямованих на розв'язання певної задачі.
\end{defn}
\subsection*{Основні властивості алгоритму}

\noindent \textbf{Скінченність} — алгоритм складається з обмеженої кількості кроків і має завершуватися після виконання.

\noindent \textbf{Визначеність (детермінованість)} — кожна команда алгоритму повинна бути чіткою й однозначною, без двозначності чи різних тлумачень.

\noindent \textbf{Масовість (загальність)} — алгоритм застосовується для широкого класу задач або різних вхідних даних, а не лише для одного конкретного випадку.

\noindent \textbf{Результативність} — виконання алгоритму повинно приводити до розв’язання задачі або отримання результату.

\noindent \textbf{Дискретність} — алгоритм складається з окремих кроків (операцій), які виконуються послідовно.


	\section{Структурне програмування}

	Структурне програмування --- це методологія розробки програм, яка ґрунтується 
	на використанні трьох базових конструкцій управління виконанням:
	\begin{itemize}
		\item слідування (лінійний алгоритм),
		\item розгалуження (умовний оператор),
		\item цикл (повторення).
	\end{itemize}
	
	\subsection*{Слідування (лінійний алгоритм)}
 
		\begin{defn}{Лінійні алгоритми — це}
		алгоритми, в яких дії виконуються послідовно одна за одною. 
	\end{defn}

	Це означає, що кожна команда лінійного алгоритму виконується строго в заданому порядку. Лінійні алгоритми можуть бути корисними для розв’язання простих задач, таких як знаходження суми двох чисел або переведення одиниць виміру.
	Наприклад, лінійний алгоритм для знаходження суми двох чисел може мати наступний вигляд:
	
	\textbf{Приклад:} обчислення суми двох чисел.
\tcbset{colback=blue!5!white, colframe=blue!75!black, arc=2mm, boxrule=0.8pt}

\begin{tcolorbox}[title=Алгоритм: Сума двох чисел]

		\begin{algorithmic}
			\STATE Ввести $a, b$
			\STATE Обчислити $s = a + b$
			\STATE Вивести $s$
		\end{algorithmic}

\end{tcolorbox}
	
	\subsection*{Розгалуження (умовний оператор)}
			\begin{defn}{Розгалуження — це}
	  алгоритмічна конструкція, яка дозволяє змінювати хід виконання програми залежно від істинності певної умови.
	\end{defn}
				\begin{defn}{	Умовний оператор — це}
	 конкретна реалізація розгалуження у мові програмування.
	Він перевіряє логічний вираз і, залежно від його значення (true або false), виконує одну з гілок коду.
	\end{defn}
	
  
	\textbf{Приклад:} перевірка парності числа.

	\tcbset{colback=blue!5!white, colframe=blue!75!black, arc=2mm, boxrule=0.8pt}
	
	\begin{tcolorbox}[title=Алгоритм: Перевірка парності]

			\begin{algorithmic}
				\STATE Ввести $n$
				\IF{$n \bmod 2 = 0$}
				\STATE Вивести ``Число парне''
				\ELSE
				\STATE Вивести ``Число непарне''
				\ENDIF
			\end{algorithmic}

	\end{tcolorbox}
	
	\subsection*{Цикл (повторення)}
		\begin{defn}{Цикл — це}
	алгоритмічна конструкція, яка забезпечує багаторазове виконання одних і тих самих дій (команд) доти, доки виконується певна умова або поки не буде досягнуто заданої кількості повторень..
	\end{defn}
	

	Цикли дозволяють:
		\begin{itemize}
		\item		автоматизувати повторювані обчислення;
		\item 	скоротити код, замінивши багаторазове дублювання;
		\item 	будувати алгоритми для обробки масивів даних, серій чисел, введення/виведення.
	\end{itemize}
	
 
	\textbf{Приклад:} обчислення суми чисел від $1$ до $n$.
	\tcbset{colback=blue!5!white, colframe=blue!75!black, arc=2mm, boxrule=0.8pt}

\begin{tcolorbox}[title=Алгоритм: Сума чисел від 1 до $n$]
	
		\begin{algorithmic}
			\STATE Ввести $n$
			\STATE $s := 0$
			\FOR{$i := 1$ \TO $n$}
			\STATE $s := s + i$
			\ENDFOR
			\STATE Вивести $s$
		\end{algorithmic}

\end{tcolorbox}
	
	\subsection*{	Теорема структурного програмування}
	У 1970‑х роках Едсгер Дейкстра та Ніклаус Вірт сформулювали принципи структурного програмування.
	
	Вони довели, що
	\begin{defn}{}
	
	 \textbf{будь-який алгоритм можна представити як композицію трьох базових блоків: слідування, розгалуження, цикл.} 
\end{defn}

	Логічна повнота конструкцій : слідування забезпечує лінійність і впорядкованість дій, розгалуження додає можливість приймати рішення (умовний вибір), цикл дозволяє повторювати дії довільну кількість разів. 	Разом ці три конструкції утворюють повний набір керуючих структур, достатній для опису будь-якого алгоритму.
	
	Усі сучасні мови програмування (Python, C++, Java) мають оператори для цих трьох конструкцій (if, for, while).
	
	Навіть складні алгоритми (наприклад, сортування, пошук, обробка даних) можна розкласти на комбінацію цих трьох структур.
	
	Підпрограми (функції, процедури) лише групують ці базові блоки, але не додають нових типів керування.


\section*{Контрольні питання}

\begin{enumerate}
	\item Що таке програмування? Які основні етапи воно включає?
	\item Дайте означення алгоритму. Які його основні властивості?
	\item Що таке структурне програмування? Які три базові конструкції воно використовує?
	\item Поясніть, що таке слідування (лінійний алгоритм). Наведіть приклад.
	\item Що таке розгалуження? Чим воно відрізняється від умовного оператора?
	\item Наведіть приклад умовного оператора у будь-якій мові програмування.
	\item Що таке цикл? Які основні види циклів ви знаєте?
	\item Наведіть приклад алгоритму з використанням циклу.
	\item Сформулюйте теорему структурного програмування. На чому ґрунтується її доведення?
	\item Чому саме три конструкції (слідування, розгалуження, цикл) вважаються достатніми для побудови будь-якого алгоритму?
\end{enumerate}

	
	
	
	
	


	% l1_content.tex
% Частина конспекту лекції "Python. Вступ та синтаксис", оформлена відповідно до вимог ОНУ.





\chapter{\parbox{0.85\textwidth}{Чому Python?  Створення та запуск файлів. Робота в інтерпретаторі.}}

\section{Історія та Глобальний вплив}

Створення Python розпочав Гвідо ван Россум (Guido van Rossum) у 1991 році під час роботи над розподіленою ОС «Амеба». Йому була потрібна розширювана мова, яка б забезпечувала підтримку системних викликів. За основу було взято мови ABC та Modula-3. Назву Гвідо обрав на честь популярного комедійного шоу BBC «Літаючий цирк Монті-Пайтона», а не на честь змії. Відтоді Python активно розвивався за підтримки організацій, у яких працював його автор. Сьогодні мова вдосконалюється ще швидше, адже над нею працює не лише команда творців, а й велика міжнародна спільнота програмістів. Втім, остаточне слово щодо напряму розвитку Python і надалі залишається за Гвідо ван Россумом.

\subsection*{Історія розвитку Python}

\begin{itemize}
	\item \textbf{Python 2.0} --- випущений у 2000 році. Серед нововведень: спискові включення, підтримка Unicode, вдосконалений збирач сміття.
	\item \textbf{Python 3.0} --- випущений у 2008 році. Основна особливість --- несумісність із попередніми версіями, що вимагало адаптації бібліотек.
	\item \textbf{Python 2.7} --- останній реліз гілки 2.x. Його підтримку офіційно завершено у 2020 році, що ознаменувало перехід до епохи Python~3.
\end{itemize}

\subsection*{Глобальний вплив та переваги}

Python став однією з найпопулярніших мов програмування у світі завдяки простоті синтаксису та універсальності. Його активно використовують такі компанії, як \textbf{Google, Netflix, Facebook, Instagram, Spotify} та багато інших.

\medskip
Основні сфери застосування:
\begin{itemize}
	\item \textbf{Наука та дослідження} --- бібліотеки NumPy, SciPy, pandas.
	\item \textbf{Машинне навчання та штучний інтелект} --- TensorFlow, PyTorch, scikit-learn.
	\item \textbf{Веб-розробка} --- фреймворки Django, Flask.
	\item \textbf{Автоматизація та DevOps} --- скрипти для адміністрування, тестування, CI/CD.
	\item \textbf{Аналіз даних та візуалізація} --- matplotlib, seaborn, Jupyter Notebooks.
\end{itemize}

\medskip
Таким чином, Python використовується практично всюди --- від навчальних проєктів до масштабних систем світових технологічних гігантів.



\section{Переваги мови або чому Python}

\subsection*{Основні переваги}

\begin{itemize}
	\item \textbf{Якість програмного забезпечення.} Зручність читання, ясність і висока якість відрізняють Python від інших мов. Єдиноманітність оформлення коду полегшує його розуміння навіть для тих, хто не брав участі у створенні програми.
	
	\item \textbf{Висока швидкість розробки.} У порівнянні з C, C++ та Java, Python значно підвищує продуктивність праці розробника. Обсяг коду зазвичай становить третину або навіть п’яту частину еквівалентного коду на C++ чи Java. Програми запускаються відразу, минаючи етапи компіляції та зв’язування.
	
	\item \textbf{Перенесення програмного коду без витрат.} Більшість програм виконується без змін на основних платформах. Перенесення з Linux у Windows часто зводиться до копіювання файлів. Python також підтримує створення графічних інтерфейсів, доступ до баз даних, веб-додатків та інших типів програм.
	
	\item \textbf{Бібліотеки підтримки.} У складі Python є велика стандартна бібліотека, що охоплює пошук тексту, роботу з мережею та інші функції. Крім того, існують численні сторонні бібліотеки: для веб-розробки, математичних обчислень, ігор тощо. Наприклад, \texttt{NumPy} --- потужний інструмент для наукових обчислень.
	
	\item \textbf{Інтеграція компонентів.} Python легко взаємодіє з іншими мовами та платформами: може викликати функції з бібліотек на C/C++, інтегруватися з Java, взаємодіяти з COM та .NET, працювати через протоколи SOAP, XML-RPC, CORBA.
	
	\item \textbf{Вільне використання.} Python можна використовувати та розповсюджувати абсолютно безкоштовно.
\end{itemize}

\subsection*{Класифікація мови}

\begin{enumerate}
	\item \textbf{Високого рівня.} Код легко читати і писати, він схожий на англійську мову.
	\item \textbf{Загального призначення.} Python вирішує широкий спектр завдань.
	\item \textbf{Інтерпретована.} Для виконання коду не потрібна компіляція.
\end{enumerate}

Python часто позиціонують як \textbf{об’єктно-орієнтовану мову сценаріїв}, що дозволяє створювати прості програми у вигляді одного файлу, використовуючи об’єкти.

\subsection*{Недоліки}

\begin{itemize}
	\item \textbf{Швидкодія.} Програми на Python можуть працювати повільніше за аналоги, написані компілюючими мовами (наприклад, C).
	\item \textbf{Байт-код.} Python перетворює вихідний код у байт-код, який інтерпретується. Це забезпечує переносимість програм, але відсутність машинного коду може знижувати продуктивність.
\end{itemize}

\section{Середовище для роботи та Інтерпретатор}


Python --- це не лише мова програмування, а й програмний пакет, що включає \textbf{інтерпретатор} --- модуль, який створює та виконує інші програми.

\subsection*{Онлайн середовища розробки}

Сьогодні є можливість писати код на Python у веб-сервісах, які дозволяють запускати програми прямо у браузері, без встановлення додаткового програмного забезпечення:
\begin{itemize}
	\item \url{https://replit.com/~}
	\item \url{https://ideone.com/}
	\item \url{https://www.online-python.com/}
\end{itemize}

\medskip
У цих сервісах достатньо створити акаунт, написати код та запустити його онлайн.

\subsection*{Командний рядок інтерпретатора}

Найпростіший спосіб писати й запускати програми на Python --- вводити інструкції безпосередньо у \textbf{командний рядок інтерпретатора}. Це дозволяє швидко перевіряти окремі команди та експериментувати з кодом.

\subsection*{Локальне встановлення Python}

Щоб працювати локально, потрібно завантажити Python з офіційного сайту:
\begin{itemize}
	\item \url{https://www.python.org/downloads/}
\end{itemize}

Після встановлення пакету Python створюється інтерпретатор та бібліотека підтримки. Також доступне інтегроване середовище розробки \textbf{IDLE}, яке встановлюється разом із Python і дозволяє писати, запускати та відлагоджувати програми у зручному графічному інтерфейсі.

\section{Виконання вказівок у Python}

\subsection*{Інтерактивний режим}

В інтерактивному режимі користувач вводить інструкції та натискає \texttt{Enter}, після чого результат відображається рядком нижче.

\begin{verbatim}
	>>> 2 + 2
	4
	>>> "Hello"
	'Hello'
\end{verbatim}

Інтерактивний режим дозволяє:
\begin{itemize}
	\item швидко перевіряти фрагменти коду;
	\item проводити експерименти без створення файлів;
	\item одразу виявляти помилки.
\end{itemize}

Приклад помилки:
\begin{verbatim}
	>>> print(y)
	Traceback (most recent call last):
	File "<pyshell#15>", line 1, in <module>
	print(y)
	NameError: name 'y' is not defined
\end{verbatim}

\subsection*{Виведення повідомлень на екран}

Для виведення інформації використовується команда \texttt{print}:

\begin{verbatim}
	>>> print("Hello World")
	Hello World
	>>> print(2 + 2)
	4
\end{verbatim}

У інтерактивному режимі можна вводити вирази без \texttt{print}, але у файлах інструкцію \texttt{print} потрібно писати обов’язково.

\subsection*{Програми у файлах}

Програму на Python можна написати у будь-якому текстовому редакторі та зберегти у файл з розширенням \texttt{.py}.  
Наприклад, створимо файл \texttt{hello.py}:

\begin{verbatim}
	print("Hello world!")
	print("Привіт світ")
	input()
\end{verbatim}

Запуск програми здійснюється подвійним клацанням по файлу.  
Такі програми виконуються у консольному вікні («чорному вікні»), де введення здійснюється з клавіатури. Якщо у програмі є помилки, у консолі буде показано повідомлення про них.

\medskip
Таким чином, Python дозволяє працювати як у інтерактивному режимі, так і створювати повноцінні програми у файлах.

\section{Перевірка коду на помилки в IDLE}

\subsection*{Запуск програми у IDLE}

\begin{enumerate}
	\item Запускаємо середовище \textbf{IDLE}.
	\item У меню \textbf{File} обираємо пункт \textbf{Open} та у вікні шукаємо файл \texttt{hello.py}.
	\item Відкриється нове вікно з текстом програми. У цьому вікні у меню \textbf{Run} вибираємо \textbf{Run Module} або натискаємо клавішу \texttt{F5}.
\end{enumerate}

\subsection*{Результати виконання}

\begin{itemize}
	\item Якщо у програмі немає помилок, введення та виведення здійснюється у вікні інтерпретатора.
	\item Повідомлення про помилки також відображаються у \textbf{IDLE}.
	\item Якщо програма працює некоректно, але не видає помилок, для пошуку проблем можна вставити додаткові інструкції \texttt{print}.
\end{itemize}

\subsection*{Приклад використання \texttt{print} для відлагодження}

\begin{verbatim}
	print("Значення змінної x:", x)
\end{verbatim}

Це допоможе зрозуміти, які значення мають змінні під час виконання програми.

\section*{Контрольні питання}

\begin{enumerate}
	\item Які основні сфери застосування Python ви знаєте?
	\item Назвіть головні переваги мови Python.
	\item Які недоліки Python були згадані у тексті?
	\item Як класифікується мова Python (за рівнем, призначенням, способом виконання)?
	\item Що таке інтерпретатор Python і яку роль він виконує?
	\item Які онлайн‑середовища можна використати для написання та запуску програм на Python?
	\item Як працювати з Python у командному рядку інтерпретатора?
	\item Які кроки потрібно виконати для локального встановлення Python?
	\item Чим відрізняється інтерактивний режим від запуску програм у файлах?
	\item Як у Python здійснюється виведення повідомлень на екран?
	\item Як перевіряти код на помилки у середовищі IDLE?
	\item Які результати виконання програми можна побачити у IDLE?

\end{enumerate}

	% l1_content.tex
% Частина конспекту лекції "Python. Вступ та синтаксис", оформлена відповідно до вимог ОНУ.





\chapter{\parbox{0.85\textwidth}{Чому Python?  Створення та запуск файлів. Робота в інтерпретаторі.}}

\section{Історія та Глобальний вплив}

Створення Python було розпочато Гвідо ван Россумом ( Guido van Rossum ) у 1991 році, коли він працював над розподіленою ОС Амеба. Йому була потрібна мова, що розширюється, який би забезпечив підтримку системних викликів. За основу були взяті ABC та Модула-3. Як назви він вибрав Python на честь комедійних серій BBC "Літаючий цирк Монті-Пітона", а зовсім не за назвою змії. З того часу Python розвивався за підтримки тих організацій, у яких Гвідо працював. Особливо активно мову вдосконалюється нині, коли з неї працює як команда творців, а й ціле співтовариство програмістів з усього світу. І все-таки останнє слово про напрям розвитку мови залишається за Гвідо ван Россумом.

\subsection{Історія розвитку Python}

\begin{itemize}
	\item \textbf{Python 2.0} --- випущений у 2000 році. Серед нововведень: спискові включення, підтримка Unicode, вдосконалений збирач сміття.
	\item \textbf{Python 3.0} --- випущений у 2008 році. Основна особливість --- несумісність із попередніми версіями, що вимагало адаптації бібліотек.
	\item \textbf{Python 2.7} --- останній реліз гілки 2.x. Його підтримку офіційно завершено у 2020 році, що ознаменувало перехід до епохи Python~3.
\end{itemize}

\subsection{Глобальний вплив та переваги}

Python став однією з найпопулярніших мов програмування у світі завдяки простоті синтаксису та універсальності. Його активно використовують такі компанії, як \textbf{Google, Netflix, Facebook, Instagram, Spotify} та багато інших.

\medskip
Основні сфери застосування:
\begin{itemize}
	\item \textbf{Наука та дослідження} --- бібліотеки \texttt{NumPy}, \texttt{SciPy}, \texttt{pandas}.
	\item \textbf{Машинне навчання та штучний інтелект} --- \texttt{TensorFlow}, \texttt{PyTorch}, \texttt{scikit-learn}.
	\item \textbf{Веб-розробка} --- фреймворки \texttt{Django}, \texttt{Flask}.
	\item \textbf{Автоматизація та DevOps} --- скрипти для адміністрування, тестування, CI/CD.
	\item \textbf{Аналіз даних та візуалізація} --- \texttt{matplotlib}, \texttt{seaborn}, \texttt{Jupyter Notebooks}.
\end{itemize}

\medskip
Таким чином, Python використовується практично всюди --- від навчальних проєктів до масштабних систем світових технологічних гігантів.



\section{Переваги мови або чому Python}

\subsection{Основні переваги}

\begin{itemize}
	\item \textbf{Якість програмного забезпечення.} Зручність читання, ясність і висока якість відрізняють Python від інших мов. Єдиноманітність оформлення коду полегшує його розуміння навіть для тих, хто не брав участі у створенні програми.
	
	\item \textbf{Висока швидкість розробки.} У порівнянні з C, C++ та Java, Python значно підвищує продуктивність праці розробника. Обсяг коду зазвичай становить третину або навіть п’яту частину еквівалентного коду на C++ чи Java. Програми запускаються відразу, минаючи етапи компіляції та зв’язування.
	
	\item \textbf{Перенесення програмного коду без витрат.} Більшість програм виконується без змін на основних платформах. Перенесення з Linux у Windows часто зводиться до копіювання файлів. Python також підтримує створення графічних інтерфейсів, доступ до баз даних, веб-додатків та інших типів програм.
	
	\item \textbf{Бібліотеки підтримки.} У складі Python є велика стандартна бібліотека, що охоплює пошук тексту, роботу з мережею та інші функції. Крім того, існують численні сторонні бібліотеки: для веб-розробки, математичних обчислень, ігор тощо. Наприклад, \texttt{NumPy} --- потужний інструмент для наукових обчислень.
	
	\item \textbf{Інтеграція компонентів.} Python легко взаємодіє з іншими мовами та платформами: може викликати функції з бібліотек на C/C++, інтегруватися з Java, взаємодіяти з COM та .NET, працювати через протоколи SOAP, XML-RPC, CORBA.
	
	\item \textbf{Вільне використання.} Python можна використовувати та розповсюджувати абсолютно безкоштовно.
\end{itemize}

\subsection{Класифікація мови}

\begin{enumerate}
	\item \textbf{Високого рівня.} Код легко читати і писати, він схожий на англійську мову.
	\item \textbf{Загального призначення.} Python вирішує широкий спектр завдань.
	\item \textbf{Інтерпретована.} Для виконання коду не потрібна компіляція.
\end{enumerate}

Python часто позиціонують як \textbf{об’єктно-орієнтовану мову сценаріїв}, що дозволяє створювати прості програми у вигляді одного файлу, використовуючи об’єкти.

\subsection{Недоліки}

\begin{itemize}
	\item \textbf{Швидкодія.} Програми на Python можуть працювати повільніше за аналоги, написані компілюючими мовами (наприклад, C).
	\item \textbf{Байт-код.} Python перетворює вихідний код у байт-код, який інтерпретується. Це забезпечує переносимість програм, але відсутність машинного коду може знижувати продуктивність.
\end{itemize}

\section{Середовище для роботи та Інтерпретатор}


Python --- це не лише мова програмування, а й програмний пакет, що включає \textbf{інтерпретатор} --- модуль, який створює та виконує інші програми.

\subsection{Онлайн середовища розробки}

Сьогодні є можливість писати код на Python у веб-сервісах, які дозволяють запускати програми прямо у браузері, без встановлення додаткового програмного забезпечення:
\begin{itemize}
	\item \url{https://replit.com/~}
	\item \url{https://ideone.com/}
	\item \url{https://www.online-python.com/}
\end{itemize}

\medskip
У цих сервісах достатньо створити акаунт, написати код та запустити його онлайн.

\subsection{Командний рядок інтерпретатора}

Найпростіший спосіб писати й запускати програми на Python --- вводити інструкції безпосередньо у \textbf{командний рядок інтерпретатора}. Це дозволяє швидко перевіряти окремі команди та експериментувати з кодом.

\subsection{Локальне встановлення Python}

Щоб працювати локально, потрібно завантажити Python з офіційного сайту:
\begin{itemize}
	\item \url{https://www.python.org/downloads/}
\end{itemize}

Після встановлення пакету Python створюється інтерпретатор та бібліотека підтримки. Також доступне інтегроване середовище розробки \textbf{IDLE}, яке встановлюється разом із Python і дозволяє писати, запускати та відлагоджувати програми у зручному графічному інтерфейсі.

\section{Виконання вказівок у Python}

\subsection{Інтерактивний режим}

В інтерактивному режимі користувач вводить інструкції та натискає \texttt{Enter}, після чого результат відображається рядком нижче.

\begin{verbatim}
	>>> 2 + 2
	4
	>>> "Hello"
	'Hello'
\end{verbatim}

Інтерактивний режим дозволяє:
\begin{itemize}
	\item швидко перевіряти фрагменти коду;
	\item проводити експерименти без створення файлів;
	\item одразу виявляти помилки.
\end{itemize}

Приклад помилки:
\begin{verbatim}
	>>> print(y)
	Traceback (most recent call last):
	File "<pyshell#15>", line 1, in <module>
	print(y)
	NameError: name 'y' is not defined
\end{verbatim}

\subsection{Виведення повідомлень на екран}

Для виведення інформації використовується команда \texttt{print}:

\begin{verbatim}
	>>> print("Hello World")
	Hello World
	>>> print(2 + 2)
	4
\end{verbatim}

У інтерактивному режимі можна вводити вирази без \texttt{print}, але у файлах інструкцію \texttt{print} потрібно писати обов’язково.

\subsection{Програми у файлах}

Програму на Python можна написати у будь-якому текстовому редакторі та зберегти у файл з розширенням \texttt{.py}.  
Наприклад, створимо файл \texttt{hello.py}:

\begin{verbatim}
	print("Hello world!")
	print("Привіт світ")
	input()
\end{verbatim}

Запуск програми здійснюється подвійним клацанням по файлу.  
Такі програми виконуються у консольному вікні («чорному вікні»), де введення здійснюється з клавіатури. Якщо у програмі є помилки, у консолі буде показано повідомлення про них.

\medskip
Таким чином, Python дозволяє працювати як у інтерактивному режимі, так і створювати повноцінні програми у файлах.

\section{Перевірка коду на помилки в IDLE}

\subsection{Запуск програми у IDLE}

\begin{enumerate}
	\item Запускаємо середовище \textbf{IDLE}.
	\item У меню \textbf{File} обираємо пункт \textbf{Open} та у вікні шукаємо файл \texttt{hello.py}.
	\item Відкриється нове вікно з текстом програми. У цьому вікні у меню \textbf{Run} вибираємо \textbf{Run Module} або натискаємо клавішу \texttt{F5}.
\end{enumerate}

\subsection{Результати виконання}

\begin{itemize}
	\item Якщо у програмі немає помилок, введення та виведення здійснюється у вікні інтерпретатора.
	\item Повідомлення про помилки також відображаються у \textbf{IDLE}.
	\item Якщо програма працює некоректно, але не видає помилок, для пошуку проблем можна вставити додаткові інструкції \texttt{print}.
\end{itemize}

\subsection{Приклад використання \texttt{print} для відлагодження}

\begin{verbatim}
	print("Значення змінної x:", x)
\end{verbatim}

Це допоможе зрозуміти, які значення мають змінні під час виконання програми.

\chapter{\parbox{0.84\textwidth}{Синтаксис Python. Типи Даних}}

\section{Змінні та Ідентифікатори}

\begin{defn}{Змінна - це}
	ім'я (ідентифікатор), яке може набувати певного значення.
\end{defn}

У Python змінні не потрібно оголошувати, значення можна надати одразу за допомогою символу \texttt{=}.

\begin{py}{Приклади присвоєння}
x = 4
s1 = "Hello"
_x = True
x2 = 23.46
summa1 = summa2 = k = 0
\end{py}

\begin{defn}{Ідентифікатор - це}
	унікальна ознака об'єкта, що дозволяє розрізняти його з іншими.
\end{defn}

\subsection{Правила побудови ідентифікаторів}

\begin{itemize}
	\item Використовуються лише латинські літери, цифри та знак підкреслення (\_).
	\item Ідентифікатор повинен починатися з літери або символу підкреслення.
	\item Регістр літер має значення (\texttt{X} $\neq$ \texttt{x}).
	\item Неприпустимі пробіли, точки, спеціальні символи та літери кирилиці.
	\item Не можна застосовувати спеціальні символи \texttt{\$}, \texttt{\%}, \texttt{@}.
	\item Заборонено використовувати зарезервовані слова Python (\texttt{if}, \texttt{for}, \texttt{while}, \texttt{class}, \texttt{def} тощо).
\end{itemize}

\section{Зарезервовані слова Python}

У Python існує набір \textbf{ключових слів}, які не можна використовувати як ідентифікатори (назви змінних, функцій, класів тощо).  
Ці слова мають спеціальне значення у синтаксисі мови.

\begin{center}
	\begin{tabular}{lllll}
		\texttt{False} & \texttt{class} & \texttt{finally} & \texttt{is} & \texttt{return} \\
		\texttt{None} & \texttt{continue} & \texttt{for} & \texttt{lambda} & \texttt{try} \\
		\texttt{True} & \texttt{def} & \texttt{from} & \texttt{nonlocal} & \texttt{while} \\
		\texttt{and} & \texttt{del} & \texttt{global} & \texttt{not} & \texttt{with} \\
		\texttt{as} & \texttt{elif} & \texttt{if} & \texttt{or} & \texttt{yield} \\
		\texttt{assert} & \texttt{else} & \texttt{import} & \texttt{pass} & \texttt{break} \\
	\end{tabular}
\end{center}

\medskip
\noindent Використання будь-якого з цих слів як назви змінної або функції призведе до помилки.

\subsection{Приклад помилки}

\begin{verbatim}
	>>> class = 5
	File "<stdin>", line 1
	class = 5
	^
	SyntaxError: invalid syntax
\end{verbatim}

\section{Оператори та Вирази}

\begin{defn}{Вираз — це}
	фрагмент мови програмування, що представляє спосіб обчислення деякого значення.
\end{defn}
\begin{defn}{Операнд — це}
	частина виразу, над яким виконуються операції.
\end{defn}
\begin{defn}{Оператор — це}
	дія, яка може бути виконана над одним або декількома операндами.
\end{defn}

\begin{py}{Приклади синтаксису}
X=4+5 % Присвоєння
Print(x)
S="Hello "+"World"
Print(s)
S="Helo"+\
"World" % Перенесення довгих операторів
\end{py}

\subsection{Правила форматування коду}
\begin{itemize}
	\item Кожен оператор завершується символом нового рядка.
	\item Відступи застосовуються для позначення різних блоків коду (тіла функцій, умовні висловлювання, цикли).
	\item Пробіли неприпустимі перед операторами, але необхідні перед вкладеними блоками (\texttt{unexpected indent}).
	\item Якщо потрібно розмістити кілька операторів в одному рядку, їх розділяють крапкою з комою (\texttt{;}).
	\item Символ \texttt{\#} означає коментар.
	\item Багаторядковий коментар укладають у три подвійні лапки в \texttt{py} файлі.
\end{itemize}

\section{Основні Типи Даних і Операції}

Вбудовані типи даних:
\begin{itemize}
	\item Числа: \texttt{1234}, \texttt{3.14}, \texttt{3+4j}.
	\item Рядки: \texttt{'spam'}, \texttt{"guido's"}.
	\item Списки: \texttt{[1, [2, 'three'], 4]}.
	\item Словники: \texttt{\{'food': 'spam', 'taste': 'yum'\}}.
	\item Кортежі: \texttt{(1, 'span', 4, 'U')}.
	\item Файли: myfile = open('eggs', ‘r')
	\item  Множини: set('abc'). {'a','b', 'с'}
	\end{itemize}
	
	
	\hspace{5cm}\textbf{Вбудовані числові типи}
	\begin{itemize}
	\item 
	цілі числа (числа без дробової частини), (3, 3465, 45, 42, -4) (з необмеженою точністю)
	\item дійсні числа (грубо кажучи, числа з десятковою точкою з плаваючою точкою) (3,15 5,68)
	\item  8-річні, починаються нуль про 0о77 = 63 0о10 = 8
	\item 16-річні, починаються нуль x 0xFF=255 0xff=255 • - двійкові, починаються нуль b 0b1111 = 15 0b10 = 2
	\item - комплексні числа з уявною частиною, (3+4j), (3+4J)(з допомогою великої та маленької) літер)
	\end{itemize}
	Символ множення перед j не вказується.
	\begin{pydef} 
>>>z1=2+3j
>>>z2=4+8J
>>>print(z1, z2)
(2+3j) (4+8j)
	\end{pydef}
	
	Комплексні числа можуть бути створені за допомогою функції complex(реч \_ частина, уявна \_ частина)\\
	
	\begin{pydef} 
>>>z3=complex(34,2)
>>>print(z3)
(34+2j)
	\end{pydef}
	
	\textbf{Операції з числовими типами}\\
	
	\begin{tabular}{|c|c|c|c|} 
	\hline
	\textbf{Оператор} & \textbf{Опис}&\textbf{Приклад}&\textbf{Результат}\\ \hline
	+ & Додавання&5+7&12\\ \hline
	-&Віднімання&55-10&45\\ \hline 
	*& Множення& 4*6 &24\\ \hline
	/& Ділення &7 / 2 &3.5\\ \hline
	//& Цілочисельне ділення& 7 // 2& 3\\ \hline
	\%& Остача від ділення& 7 \%  3 &1\\ \hline
	**& Піднесення до степеня& 3 ** 4 &81\\ \hline
	
	\end{tabular}
	\\
	
	Результатом операції буде новий об'єкт. \\
	 2 \bullet \bullet  100 \# 2 у степені 100\\ 
	1267650600228229401496703205376 
	
	Зверніть увагу на результат останньої операції: у Python 3 цілі числа автоматично забезпечують \textbf{необмежену} точність. Наприклад, ви можете спробувати обчислити 2 ступенем 1 000 000 (але навряд чи варто це робити, оскільки на екран буде виведено число довжиною понад 300 000 знаків, що може зайняти тривалий час!).\\
	
	\textbf{ > len(str(2 ** 1000000))} 
	
	Скільки цифр у дійсно \textsc{великому} числі\\
	301030
		
	\section{Приклади. Виділення цифр числа}
	
	\subsection{Двоцифрове число} Нехай маємо число \texttt{n = 47}.  	Щоб виділити окремі цифри:
	
	\begin{py}{Двоцифрове число}
n = 47
d1 = n // 10   # перша цифра (десятки)
d2 = n % 10    # друга цифра (одиниці)
print("Десятки =", d1, "Одиниці =", d2)
	\end{py}
	
	Результат:
	\begin{verbatim}
Десятки = 4 Одиниці = 7
	\end{verbatim}
	
	\subsection{Трицифрове число}
	
	Нехай маємо число \texttt{m = 583}.  
	Щоб виділити окремі цифри:
	
	\begin{py}{Трицифрове число}
m = 583
sot = m // 100          # сотні
des = (m // 10) % 10    # десятки
od = m % 10             # одиниці
print("Сотні =", sot, "Десятки =", des, "Одиниці =", od)
	\end{py}
	
	Результат:
	\begin{verbatim}
Сотні = 5 Десятки = 8 Одиниці = 3
	\end{verbatim}
	
	
	
	
	\emph{\textbf{Пріоритет операторів - порядок виконання  }}\href{https://pythonguide.rozh2sch.org.ua/#operators-priority}{тут}
	
	\section{Вбудовані функції для роботи з числами}
	
	Python надає низку \textbf{вбудованих функцій}, які можна застосовувати до числових типів.
	
	\subsection{Функція \texttt{abs(a)}}
	Повертає модуль числа.
	\begin{py}{Приклад}
x = -12
print(abs(x))   # 12
	\end{py}
	
	\subsection{Функція \texttt{divmod(a, b)}}
	Виконує ділення \texttt{a} на \texttt{b} та повертає пару чисел: частку і залишок.
	\begin{py}{Приклад}
x = 23
y = 7
print(divmod(x, y))   # (3, 2)
	\end{py}
	
	\subsection{Функція \texttt{hex(x)}}
	Конвертує ціле число у рядок із шістнадцяткових цифр.
	\begin{py}{Приклад}
x = 15
print(hex(x))   # '0xf'
	\end{py}
	
	\subsection{Функція \texttt{oct(x)}}
	Конвертує ціле число у рядок із вісімкових цифр.
	\begin{py}{Приклад}
x = 15
print(oct(x))   # '0o17'
	\end{py}
	
	\subsection{Функція \texttt{round(x)}}
	Округлює число до найближчого цілого.
	\begin{py}{Приклад}
print(round(3.0254))   # 3
	\end{py}
	
	\subsection{Функція \texttt{round(x, n)}}
	Округлює число \texttt{x} з точністю до \texttt{n} знаків після коми.
	\begin{py}{Приклад}
print(round(3.0254, 2))   # 3.03
	\end{py}
	
	\chapter{Умовний оператор if}
\justifying
За замовчуванням хід виконання програми є лінійним: інструкції виконуються послідовно — рядок за рядком, від першого до останнього. Послідовне виконання є базою, але реальні алгоритми вимагають гнучкості та адаптивності.

У більшості задач програма повинна приймати рішення залежно від вхідних даних чи поточного стану обчислень. Для цього використовуються розгалуження (або умовні переходи).

\begin{defn}{Оператор if в Python — це}
	умовний оператор, який дозволяє програмі приймати рішення та виконувати певний блок коду лише тоді, коли виконується задана умова
\end{defn}

В основі прийняття рішень у програмі лежать оператори порівняння (відношення). Вони дозволяють зіставляти значення виразів і повертають результатом булеве значення — \textbf{ True (істина)} або \textbf{False (хибність)}.
\section {Оператори порівнювання}
Python підтримує базові математичні операції відношення:

\begin{pydef}
Дорівнює:                a == b
Не дорівнює:             a != b
Менше ніж:               a < b
Менше або дорівнює:     a <= b
Більше ніж:              a > b
Більше або дорівнює:    a >= b
\end{pydef}

Зверніть увагу:

Оператор перевірки на \textbf{ рівність} позначено двома знаками \textbf{==}, оскільки один знак = в Python зарезервовано для операції \textbf{присвоєння}.

На відміну від багатьох інших мов програмування, Python дозволяє використовувати ланцюгові порівняння (подвійні нерівності), що звично для математиків, наприклад:\textbf{ a < x <= b} замість запису через логічне «і».

Отримані булеві вирази є фундаментом для керування потоком виконання програми: вони використовуються в умовних конструкціях (if, elif, else) та циклах (while).

\section{Синтаксичні форми умовного оператора}

Залежно від кількості напрямків (гілок) виконання програми, в Python застосовують три основні форми розгалуження:

\vspace{1em}

\noindent\textbf{1. Неповна форма (коротке розгалуження)}

Використовується, коли певну дію потрібно виконати \textbf{тільки} за умови, що вираз-умова є істинним (\texttt{True}). Якщо умова хибна (\texttt{False}), блок коду просто пропускається.

\begin{pydef} 
Неповна форма
if вираз_умова1:
    оператор1
    оператор2
\end{pydef}


\vspace{0.5em}

\noindent\textbf{2. Повна форма (двохгілкове розгалуження)}

Застосовується, коли є два альтернативних шляхи: один виконується при істинності умови (\texttt{True}), а інший (блок \texttt{else}) — якщо умова хибна (\texttt{False}).

\begin{pydef}
Повна форма
if вираз_умова1:
    оператор1
    оператор2
else:
    Блок_операторів_3
\end{pydef}

\vspace{0.5em}

\noindent\textbf{3. Багатогілкова форма (багатоальтернативне розгалуження)}

Використовується для перевірки серії послідовних умов за допомогою конструкцій \texttt{elif} (скорочення від \textit{else if}). Умови перевіряються зверху вниз: виконується блок першої умови, яка виявиться істинною, після чого весь оператор \texttt{if} завершує роботу.

\begin{pydef}
Багатогілкова форма
if вираз_умова1:
    Блок_операторів_1
elif вираз_умова2:
    Блок_операторів_2
elif вираз_умова_n:
    Блок_операторів_n
else:
   Блок_операторів_n+1
\end{pydef}

\vspace{0.5em}

\textbf{Важливі синтаксичні особливості Python:}
\begin{itemize}
	\item \textbf{Двокрапка (\texttt{:}):} є обов'язковою в кінці кожного рядка з \texttt{if}, \texttt{elif} та \texttt{else}. Вона позначає початок відповідного блоку коду.
\end{itemize}

\begin{py}{Приклад 1. Користувач вводить число та отримує відповідь яке воно позитивне негативне або нуль.} 
x=int(input("Введіть x "))
if x>0:
    print("x позитивне")
else:
    print("x негативне або нуль")
    input()
\end{py}

\begin{py}{Приклад 2. Користувач вводить два числа та отримує відповідь яке більше або дорівнюють. }
A=input()
B=input()
if A>B: print("А більше") #можна писати у тому ж рядку 1 оп-р
elif A==B: print("А дорівнює В")
else: print("А менше В")
\end{py}

\begin{py}{Приклад 3. Користувач отримує відповідь чи оплачений чек (ціна записана у коді, або можна ввести з клавіатури). (звернути увагу на імена змінних}
print ("Привіт")
price1 = 50 #price1=int(input()) 
price2 = 32 #price2=int(input())
if price1+ price2 > 99 :
    print ("Сума не достатня")
else:
    print ("Чек оплачений")
    print ("До побачення")
\end{py}

\begin{py}{Приклад 4. Вивести на екран місяці відповідної попи року}
s=input("Введіть пору року")
if s=="зима":
    print("грудень")
    print("січень")
    print("лютий")
elif s=="весна":
    print("березень")
    print("квітень")
    print("травень")
elif s=="літо":
    print("червень")
    print("липень")
    print("серпень")
else:
    print("вересень")
    print("жовтень ")
    print("листопад")
print ("До побачення")
\end{py}
\section*{Контрольні питання}

\begin{enumerate}
	\item Що таке умовний оператор \texttt{if} у Python? Для чого він використовується?
	\item Які оператори порівняння підтримує Python? Наведіть приклади.
	\item Чим відрізняється оператор присвоєння \texttt{=} від оператора перевірки на рівність \texttt{==}?
	\item Що таке ланцюгові порівняння у Python? Наведіть приклад.
	\item Які основні синтаксичні форми умовного оператора існують у Python?
	\item У чому полягає відмінність між неповною та повною формою умовного оператора?
	\item Як працює багатогілкова форма умовного оператора з \texttt{elif}?
	\item Яка роль двокрапки \texttt{:} у синтаксисі умовних операторів?
	\item Що відбувається, якщо умова у операторі \texttt{if} є хибною?
	\item Як можна записати умовний оператор у одному рядку?
	
\end{enumerate}
\section*{Практичні завдання}
\begin{enumerate}
	\item Реалізуйте програму, яка за двома введеними числами перевіряє, чи можна знайти частку від ділення першого на друге.
	\item Напишіть програму, яка за введеними координатами точки возначає чверть.
	\item Складіть програму, яка визначає найбільше з трьох введених чисел.
	\item Реалізуйте програму, яка перевіряє, чи введене число є парним або непарним.
	\item Напишіть програму, яка визначає, чи введене число кратне 10.
	\item Складіть програму, яка за введеним віком користувача виводить повідомлення: «дитина», «підліток», «дорослий», «пенсіонер».
	\item Реалізуйте програму, яка перевіряє, чи введений рік є високосним.
	\item Напишіть програму, яка за введеною оцінкою (від 1 до 12) виводить текстове пояснення («незадовільно», «задовільно», «добре», «відмінно»).
\end{enumerate}

	

\chapter{\parbox{0.84\textwidth}{Логічні оператори (and, or, not)}}

Прості умовні вирази (операції відношення) можна об'єднувати у складні \textbf{складені умови} за допомогою \textbf{логічних операторів} (булевих зв'язок). В Python використовуються три базові логічні оператори: \textbf{\texttt{and}, \texttt{or} та \texttt{not}}.

\vspace{0.5em}

\noindent\textbf{1. Кон'юнкція (\texttt{and} — логічне «І»)}

Оператор \textbf{\texttt{and}} повертає \textbf{\texttt{True}} тоді й лише тоді, коли \textbf{обидва} операнди є істинними. Якщо хоча б один із виразів хибний (\textbf{\texttt{False}}), увесь складений вираз набуває значення \textbf{\texttt{False}}.

Пріоритет оператора \textbf{ \texttt{and}} нижчий за пріоритет операторів порівняння (\texttt{>}, \texttt{<}, \texttt{==} тощо), тому у виразах на зразок:
\begin{py}{Перевірка діапазону}
if length > 8 and length < 20:
    print("Довжина відповідає вимогам")
\end{py}
\noindent дужки навколо окремих порівнянь писати не обов'язково.

\vspace{0.5em}

\noindent\textbf{2. Диз'юнкція (\texttt{or} — логічне «АБО»)}

Оператор \textbf{\texttt{or}} повертає \textbf{\texttt{True}}, якщо \textbf{хоча б один} із операндів є істинним (або обидва одночасно). Вираз набуває значення \textbf{\texttt{False}} лише тоді, коли всі складові частини є хибними.

\begin{py}{Приклад диз'юнкції}
if age < 18 or age >= 60:
    print("Пільгова категорія")
\end{py}

\vspace{0.5em}

\noindent\textbf{3. Заперечення (\texttt{not} — логічне «НЕ»)}

Унарний оператор \textbf{\texttt{not}} інвертує (змінює на протилежне) значення булевого виразу: \textbf{\texttt{True}} перетворюється на \textbf{\texttt{False}}, а \textbf{\texttt{False}} — на \textbf{\texttt{True}}.

\begin{py}{Приклад заперечення}
if not is_valid:
    print("Помилка введення даних")
\end{py}

\vspace{0.5em}

\begin{table}[htbp]
	\centering
	\caption{Таблиця істинності для логічних операцій}
	\begin{tabular}{|c|c|c|c|c|}
		\hline
		\multicolumn{2}{|c|}{Значення операндів} & \multicolumn{3}{c|}{Результат операції} \\ \hline
		X & Y & not X & X and Y & X or Y \\ \hline
		False & False & True  & False   & False  \\ \hline
		False & True  & True  & False   & True   \\ \hline
		True  & False & False & False   & True   \\ \hline
		True  & True  & False & True    & True   \\ \hline
	\end{tabular}
\end{table}

\vspace{0.5em}

\textbf{Особливості Python:}
\begin{itemize}
	\item \textbf{Пріоритет виконання:} Найвищий пріоритет має \texttt{not}, далі йде \texttt{and}, і найнижчий — \texttt{or}. Для зміни порядку обчислення використовують круглі дужки \texttt{()}.
	\item \textbf{Ліниві обчислення (Short-circuit evaluation):} Python обчислює логічні вирази зліва направо і зупиняється, як тільки результат стає однозначним (наприклад, якщо у виразі з \texttt{and} перший операнд \texttt{False}, другий операнд навіть не обчислюється).
\end{itemize}
\chapter{\parbox{0.75\textwidth}{Генерація випадкових чисел: модуль random}}

При розв'язанні багатьох математичних задач, моделюванні ймовірнісних процесів (метод Монте-Карло) або створенні ігор виникає потреба у використанні **випадкових чисел**. У Python для цього призначений вбудований модуль \texttt{random}.

Щоб отримати доступ до функцій модуля, його необхідно імпортувати на початку програми:
\begin{py}{Імпорт модуля}
import random
\end{py}

\section*{Основні функції модуля random}

\noindent\textbf{1. Дійсні числа: \texttt{random.random()} та \texttt{random.uniform()}}

\begin{itemize}
	\item \texttt{random.random()} — повертає випадкове дійсне число $x$ у напіввідкритому інтервалі $[0.0, 1.0)$.
	\item \texttt{random.uniform(a, b)} — повертає випадкове дійсне число $x$ у відрізку $[a, b]$.
\end{itemize}

\begin{py}{Генерація дійсних чисел}
import random
# Випадкове число від 0.0 до 1.0
p = random.random()
print("Імовірність",p)

# Випадкова температура в межах [-5.0, 15.0]
temp = random.uniform(-5.0, 15.0)
print("Температура:", temp)
\end{py}

\vspace{0.5em}

\noindent\textbf{2. Цілі числа: \texttt{random.randint()} та \texttt{random.randrange()}}

\begin{itemize}
	\item \texttt{random.randint(a, b)} — повертає випадкове ціле число $n \in [a, b]$ (включаючи обидва кінці $a$ та $b$).
	\item \texttt{random.randrange(start, stop, step)} — повертає випадковий елемент із послідовності \texttt{range(start, stop, step)}.
\end{itemize}

\begin{py}{Генерація цілих чисел}
import random
# Моделювання кидка грального кубика (від 1 до 6)
dice = random.randint(1, 6)
print("Випало число:", dice)

# Випадкове парне число від 10 до 98
even_num = random.randrange(10, 100, 2)
print("Парне число:",even_num)
\end{py}

\vspace{0.5em}


\section*{Відтворюваність результатів: \texttt{random.seed()}}

Оскільки Python генерує **псевдовипадкові числа** за допомогою алгоритму (Вихр Мерсенна), послідовність обчислень залежить від початкового значення (значення ініціалізації — \textit{seed}).

Функція \texttt{random.seed(a)} фіксує початковий стан генератора. Це є критично важливим у наукових та математичних обчисленнях для **відтворюваності результатів експерименту**.

\begin{py}{Фіксація генератора}
import random

random.seed(42)  # Задаємо фіксоване зерно
print(random.randint(1, 100))  # Завжди повертатиме один і той самий результат при повторних запусках
\end{py}



	\chapter{\parbox{0.75\textwidth}{Оператор циклу \texttt{for} та функція \texttt{range()}}}
	
Оператори циклів використовують у випадках, коли в програмі необхідно багаторазово виконати певну послідовність команд. У мові Python основним циклом із заздалегідь відомою кількістю повторень є цикл \texttt{for}.

\vspace{0.3cm}

\begin{pydef}
Синтаксис циклу for
for параметр in range(start, stop, step):
    блок_операторів
\end{pydef}

\vspace{0.3cm}

\section*{1. Генератор послідовностей \texttt{range()}}

Найчастіше для формування послідовностей чисел у циклі \texttt{for} використовується вбудована функція \texttt{range()}. Вона може приймати від 1 до 3 аргументів:

\begin{itemize}
	\item \texttt{range(stop)} "--- генерує числа від $0$ до $stop - 1$ із кроком $1$.
	\item \texttt{range(start, stop)} "--- генерує числа від $start$ до $stop - 1$ із кроком $1$.
	\item \texttt{range(start, stop, step)} "--- генерує числа від $start$ до $stop - 1$ із кроком $step$.
\end{itemize}

\begin{table}[h!]
	\centering
	\renewcommand{\arraystretch}{1.3}
	\begin{tabular}{|l|l|l|}
		\hline
		\textbf{Виклик \texttt{range()}} & \textbf{Згенерована послідовність} & \textbf{Примітка} \\ \hline
		\texttt{range(5)} & \texttt{0, 1, 2, 3, 4} & Від 0 до 4 \\ \hline
		\texttt{range(1, 6)} & \texttt{1, 2, 3, 4, 5} & Від 1 до 5 \\ \hline
		\texttt{range(2, 10, 2)} & \texttt{2, 4, 6, 8} & Парні числа від 2 до 8 \\ \hline
		\texttt{range(10, 0, -1)} & \texttt{10, 9, 8, 7, 6, 5, 4, 3, 2, 1} & Зворотний відлік \\ \hline
	\end{tabular}
	\caption{Приклади роботи функції \texttt{range()}}
\end{table}

\begin{py}{Якщо передати в range() лише одне число, за замовчуванням відлік починається з 0, а верхня межа не включається}
for i in range(5):
    print(i)
\end{py}
Результат виведення : 0, 1, 2, 3, 4


\begin{py}{range(start, stop) починає з числа start (включно) і зупиняється перед stop (не включаючи його).}
for i in range(1,5):
    print(i)
\end{py}
Результат виведення: 1, 2, 3, 4

\begin{py}{Третій параметр — це крок (step). range(1, 10, 2) починає з 1 і додає 2 на кожному кроці, генеруючи непарні числа до 10}
for i in range(1,10,2):
    print(i)
\end{py}
Результат виведення: 1, 3, 5, 7, 9

\begin{py}{Якщо крок від'ємний (-1), значення будуть зменшуватися. Цикл починається з 10 і зупиняється перед 0.}	
for i in range(10,0,-1):
    print(i)
\end{py}
Результат виведення: 10, 9, 8, 7, 6, 5, 4, 3, 2, 1

\section*{2. Базові алгоритми з циклом \texttt{for}}

\begin{py}{Приклад 1. Обчислення суми натуральних чисел від 1 до n}
n = int(input("Введіть n "))
s = 0
for i in range(1, n + 1):
    s += i

print(f"Сума чисел від 1 до {n} дорівнює: {s}")
\end{py}

input(...) запитує у користувача число та приймає його як рядок (текст).

int(...) перетворює введений текст у ціле число та зберігає його в змінну $n$.

$s = 0$ cтворюється змінна-накопичувач $s$ із початковим значенням 0. У ній буде зберігатися підсумкова сума.

for i in range(1, n + 1): цикл for перебирає числа по порядку.

range(1, n + 1) генерує послідовність від 1 до $n$ включно (верхня межа $n + 1$ не включається в діапазон, тому цикл зупиняється саме на n).

s += i  на кожній ітерації циклу поточне число i додається до суми s (це скорочений запис для $s = s + i).$

print(f"Сума чисел від 1 до {n} дорівнює: {s}")

За допомогою $f$-рядка на екран виводиться підсумковий результат.
\begin{py}{Приклад 2. Обчислення факторіала n!}
n = int(input("Введіть n "))
fact = 1

for i in range(1, n + 1):
    fact *= i

print(f"{n}! = {fact}")
\end{py}

\begin{py}{Приклад 3. Табулювання функції y = sin(x) на відрізку [a, b]}
import math

a = float(input("Введіть початок відрізка a "))
b = float(input("Введіть кінець відрізка b "))
n = 10
h = (b - a) / n

print(" x | sin(x)")
print("-------------")
for i in range(n + 1):
    x = a + i * h
    y = math.sin(x)
    print(f"{x:6.2f} | {y:6.4f}")
\end{py}


Форматоване виведення (виділений рядок):

$\{x:6.2f\}$ — виводить значення $x$: занімає 6 символів у ширину та округлюється до 2 знаків після коми.| — виводить розділювальну риску між стовпчиками.

$\{y:6.4f\}$ — виводить значення $y$ ($\sin(x)$): шириною в 6 символів та округлене до 4 знаків після коми.

\vspace{0.5cm}

\section*{3. Аналіз дільників числа та перевірка на простоту}

\begin{tcolorbox}[colback=pyblue!5!white,colframe=pyblue,title=\textbf{Математична довідка},fonttitle=\sffamily\bfseries]
	\textbf{Просте число} "--- це натуральне число, більше за 1, яке має точно два різні натуральні дільники: 1 та саме це число. Якщо число $n$ має дільник, відмінний від 1 та $n$, то найменший такий дільник не перевищує $\lfloor\sqrt{n}\rfloor$.
\end{tcolorbox}

\begin{py}{Приклад 4. Виведення всіх дільників числа}
n = int(input("Введіть натуральне число n "))
count=0
for i in range(1, n + 1):
    if n % i == 0:
        print(i)
    count=count+1
print(f"Кількість дільників: {count}")
\end{py}

Пошук дільників: Програма перебирає всі цілі числа $i$ від $1$ до $n$ включно і перевіряє умову n % i == 0 (чи ділиться $n$ на $i$ без залишку).

Виведення та лічильник: Якщо число $i$ є дільником:воно одразу виводиться на екран (print(i)); лічильник count збільшується на $1$.  Наприкінці програма виводить загальну кількість знайдених дільників.



\section*{4. Обчислення нескінченних математичних рядів}

\subsection*{4.1. Сума геометричного ряду}
Обчислимо суму $S = \sum_{n=0}^{N} x^n = 1 + x + x^2 + \dots + x^N$ ($|x| < 1$).

\begin{py}{Приклад 5. Сума геометричного ряду}
x = float(input("Введіть x (|x| < 1) "))
n = int(input("Введіть кількість членів N "))

s = 0.0
term = 1.0

for i in range(n):
    s += term
    term *= x

print(f"Сума ряду: {s:.6f}")
print(f"Точна сума 1/(1-x): {1 / (1 - x):.6f}")
\end{py}

\subsection*{4.2. Обчислення експоненти $e^x$ через ряд Тейлора}

\begin{py}{Приклад 6. Обчислення e^x через ряд Тейлора}
import math

x = float(input("Введіть x "))
n = int(input("Кількість членів ряду "))

s = 0.0
term = 1.0

for i in range(1, n + 1):
    s += term
    term *= x / i

print(f"Наближене значення e^{x} = {s:.8f}")
print(f"Точне значення (math.exp): {math.exp(x):.8f}")
\end{py}
$f$"..." ($f$-рядок) — дозволяє вставляти значення змінних безпосередньо всередину тексту за допомогою фігурних дужок {}.

$\{s:.8f\}$ — бере значення змінної $s$ (суму ряду) та форматує його як число з плаваючою крапкою, залишаючи рівно 8 знаків після коми.

Цей рядок показує наближене значення експоненти $e^x$ , вираховане через суму перших $n$ членів ряду Тейлора.
\section*{5. Управління циклом: \texttt{break}, \texttt{continue} та \texttt{else}}

\begin{itemize}
	\item \texttt{break} "--- достроково припиняє виконання циклу.
	\item \texttt{continue} "--- пропускає залишок тіла циклу та переходить до наступної ітерації.
	\item \texttt{else} "--- блок коду, який виконується після завершення циклу, якщо вихід з нього \textbf{не} відбувся за допомогою \texttt{break}.
\end{itemize}

\begin{py}{Приклад 7. Пошук першого від'ємного числа з використанням break та else}
for i in range(5):
    val = int(input("Введіть число "))
    if val < 0:
        print(f"Знайдено перше від'ємне число: {val}")
        break
else:
    print("Усі введені числа є невід'ємними.")
\end{py}

\vspace{0.3cm}

\section*{6. Оператори інкременту та особливості \texttt{float}}

У Python відсутні оператори \texttt{++} та \texttt{--}. Замість них використовуються комбіновані оператори присвоєння:
\begin{itemize}
	\item \texttt{x += 1} замість \texttt{x = x + 1}
	\item \texttt{x -= 1} замість \texttt{x = x - 1}
	\item \texttt{x *= 2} замість \texttt{x = x * 2}
\end{itemize}

\begin{tcolorbox}[colback=yellow!10!white,colframe=orange,title=\textbf{Обережно з дійсними числами (float)!},fonttitle=\sffamily\bfseries]
	Числа з плаваючою крапкою зберігаються в пам'яті комп'ютера у двійковій системі з обмеженою точністю (53 біти). Через це обчислення можуть містити незначну похибку:
\begin{py}{Особливості float}
>>> 0.1 + 0.1 + 0.1
0.30000000000000004
\end{py}
	Для порівняння дійсних чисел завжди використовуйте функцію \texttt{math.isclose(a, b)} або порівняння з точністю $\varepsilon$: \texttt{abs(a - b) < 1e-9}.
\end{tcolorbox}

\vspace{0.3cm}

\
\end{enumerate}

	\chapter{Оператор циклу \texttt{while}}

Оператори циклів використовують у випадках, коли в програмі необхідно багаторазово виконати певну послідовність команд. У мові Python основним циклом із умовою повторення (циклом із передумовою) є цикл \texttt{while}.

\vspace{0.3cm}

\begin{pydef}
Синтаксис циклу while
while вираз_умови:
    блок_операторів
\end{pydef}

\vspace{0.3cm}

Блок операторів циклу \texttt{while} повторюється до того часу, поки результатом логічного висловлювання умови повтору є значення \texttt{True}. Вираз умови обчислюється перед кожним повтором циклу, і вихід з циклу (тобто перехід до оператора, наступного за оператором \texttt{while}) здійснюється в тому випадку, якщо вираз умови має значення \texttt{False}.

\begin{py}{Приклад 1. Виведення значень лічильника від 0 до 4}
i = 0
while i < 5:
    print("значення лічильника =", i)
    i += 1  # i = i + 1

print("end")
\end{py}

\noindent\textbf{Покроковий аналіз виконання:}

Задається початкове значення змінної $i = 0$.

 Перевіряється умова $i < 5$. Оскільки вона виконується ($\text{\texttt{True}}$), відбувається перехід до операторів всередині циклу (виводимо значення $i$ на екран та збільшуємо на 1, тобто $i = 0 + 1 = 1$). 
 
 Знову повертаємось до перевірки умови на початку $i < 5$. Вона виконується $\rightarrow$ виконуються оператори виводу на екран та збільшення $i$, і так до тих пір, поки значення $i$ не стане дорівнювати 5. Тоді оператори всередині циклу вже не будуть виконуватись, виконання перейде до оператора \texttt{print("end")}.
 \subsection*{Різниця між циклами for та while}
 
 \begin{itemize}
 	\item \textbf{Цикл for} використовується тоді, коли кількість повторень відома наперед.  
 	Він перебирає елементи послідовності (наприклад, \texttt{range()}, список, рядок) і виконує тіло циклу для кожного елемента.
 	
 	\item \textbf{Цикл while} застосовується тоді, коли кількість повторень невідома заздалегідь.  
 	Він виконує тіло циклу доти, доки умова залишається істинною (\texttt{True}).  
 	Якщо умова стає хибною (\texttt{False}), цикл завершується.
 \end{itemize}
 
 

\section{ Керування циклом: \texttt{break}, \texttt{continue}, \texttt{pass} та блок \texttt{else}}

\begin{pydef}
Структура керування циклом while
while вираз_умови_test:
    блок_операторів_1
    if умова_1:
        break     # вийти з циклу, пропустивши else
    if умова_2:
        continue  # перейти на початок циклу до 
                  #виразу_умови_test
    блок_операторів_2
else:
    оператори_else # виконуються, якщо вираз_умови 
                   # став False (без break)
\end{pydef}

\begin{itemize}
	\item \textbf{\texttt{break}} "--- вихід з циклу із пропуском блоку \texttt{else}.
	\item \textbf{\texttt{continue}} "--- перехід на початок циклу до перевірки умови.
	\item \textbf{\texttt{pass}} "--- порожня інструкція. Не виконує жодних дій і використовується у випадках, коли синтаксис мови потребує наявності інструкції, але поки що жодних корисних дій виконати не потрібно:
\begin{minted}[fontsize=\small,frame=single,bgcolor=MintCream]{python}
if True:
    pass  # інструкції будуть написані пізніше

while True:
    pass  # порожній нескінченний цикл
\end{minted}
\end{itemize}

\section{ Типові застосування циклу \texttt{while}}

\begin{py}{Приклад 2. Інтерактивна перевірка умов (перевірка стану)}
x = input("Двері відчинені? ")
while x.lower() != "так":
    print("стукаємо")
    x = input("Двері відчинені? ")
print("Заходимо")
\end{py}

\textit{Примітка:} Блок операторів у \texttt{while} може не виконатись жодного разу, якщо вираз-умова від самого початку має значення \texttt{False}.

\begin{py}{Приклад 3. Табулювання функції y = sin(x) до першого від'ємного значення}
import math

x = 0.0
while math.sin(x) >= 0:
    print(f"sin({x:.1f}) = {math.sin(x):.4f}")
    x += 0.1
\end{py}

\begin{py}{Приклад 4. Запит на введення додатнього числа з валідацією}
x = int(input("Введіть додатнє x = "))
while x <= 0:
    x = int(input("Введіть додатнє x = "))
print("Дякую")
\end{py}
Програма буде виводити на екран повідомлення та перевіряти, яке число ввід користувач. Якщо $x<=0$, то повідомлення з запрошенням вводу буде з'являтися знов та знов. Якщо додатнє, то цикл завершиться, бо умова $x<=0$ перестане виконуватись, та виконаэться рядок коду, що йде одрвзу після циклу while. 

\begin{py}{Приклад 5. Обчислення суми введених чисел до введення 0 (стоп-сигнал)}
s = 0.0
x = float(input())
while x != 0:
    s += x
    x = float(input())
print("Сума =", s)
\end{py}

Ініціалізується змінна $s$, яка буде накопичувати суму введених чисел.

Користувач вводить число з клавіатури, яке перетворюється у тип float (дійсне число).

Запускається цикл, який працює доти, доки введене число не дорівнює нулю.

До змінної s додається введене число x.

Користувач вводить наступне число. Якщо це знову не нуль — цикл повторюється.

Коли користувач вводить 0, цикл завершується, і програма виводить загальну суму всіх введених чисел.

\begin{py}{Приклад 6. Обчислення суми нескінченного ряду (розклад exp(x)) з точністю eps}
# Розклад e^x = 1 + x + x^2/2! + ... + x^n/n! + ...
import math
# Початкові дані та параметри
x = 1.0       # Значення x, для якого обчислюємо e^x
eps = 1e-9    # Точність обчислень (10^-9)
s = 0.0       # Змінна для накопичення суми ряду
term = 1.0    # Перший доданок ряду (при n = 0, x^0 / 0! = 1)
i = 1         # Лічильник ітерацій та порядковий номер 
#наступного додатка

# Цикл виконується, поки модуль поточного доданка більший за 
#задану точність eps
while abs(term) > eps:
    s += term           # Додаємо поточний доданок до  суми
    term = term * x / i # Обчислюємо наступний доданок
    i += 1              # Збільшуємо номер ітерації на 1
print(f"Сума s = {s}")
print(f"Ітерацій: {i - 1}")
\end{py}

\textit{Примітка:} Будь-який цикл \texttt{for} можна записати за допомогою \texttt{while}, але зворотне твердження взагалі кажучи невірно, оскільки \texttt{while} призначений для циклів з невідомою заздалегідь кількістю повторень.

\section{Проблема «зациклювання»}

Іноді при неуважному використанні оператора \texttt{while} виникає так зване \textbf{зациклювання} "--- нескінченне повторення тіла циклу. Це може статися з двох основних причин:

\begin{enumerate}[label=\arabic*.]
	\item Змінні, що є частиною висловлювання умови циклу, не змінюють усередині циклу своїх значень або приймають некоректні значення, у результаті чого вираз умови завжди має значення \texttt{True}.
	\item Некоректним є сам вираз умови циклу (наприклад, порівняння дійсних чисел типу \texttt{float} на строгу рівність).
\end{enumerate}

\begin{py}{Приклад зациклювання 1: Нескінченний цикл при перериванні Ctrl+C}
while True:
    print("Натисніть Ctrl+C для зупинки")
\end{py}
while True: — створює нескінченний цикл. Оскільки умова True (істина) завжди залишається True і ніколи не змінюється на False, цикл буде виконуватися вічно.

print("Натисніть Ctrl+C для зупинки") — тіло циклу, яке безперервно виводить це повідомлення в консоль/термінал.

Зациклювання виникає тому, що у коді немає жодної умови виходу (наприклад, оператора break або зміни логічного виразу), тому програма самостійно зупинитися не може.

Щоб примусово перервати виконання нескінченного циклу в терміналі чи консолі, використовують комбінацію клавіш Ctrl + C. Вона надсилає програмі сигнал переривання (KeyboardInterrupt), що викликає зупинку виконання коду.
\begin{py}{Приклад зациклювання 2: Небезпечне порівняння дійсних чисел (float)}
# Некоректний код 
x = 0.0
while x != 1.0:# умова не перестане виконуватись ніколи
    print(f"x = {x:.1f}")
    x += 0.3

# Правильний варіант:
x = 0.0
while x < 1.0:
    print(f"x = {x:.1f}")
    x += 0.3
\end{py}

\section{Складніші приклади та алгоритми}

\begin{py}{Приклад 7. Середнє арифметичне простих чисел до введення 0}
x = int(input())  # Зчитуємо перше ціле число з клавіатури
k = 0             # Лічильник знайдених простих чисел
S = 0             # Змінна для збереження суми простих чисел

# Цикл триває, поки користувач не введе 0
# ( 0 - ознака кінця введення)
while x != 0:
    is_prime = True  # Прапорець: припускаємо,
                     #  що число x є простим
    if x <= 1:
        is_prime = False  # Числа <= 1 не є простими
    else:
        # Натсупний цикл перевіряє наявність дільників 
        # від 2 до  кореня з x включно
        for i in range(2, int(x**0.5) + 1):
            if x % i == 0: # Якщо x ділиться на i без остачі 
                is_prime = False  # то воно не просте, і далі
                break             # можна не перевіряти

    # Якщо число виявилося простим
    if is_prime:
        k += 1  # Збільшуємо кількість простих чисел на 1
        S += x  # Додаємо число до загальної суми

    x = int(input())  # Зчитуємо наступне число для наступної 
                      # ітерації циклу

# Підбиття підсумків після виходу з циклу
if k == 0: # Якщо жодного простого числа не було введено
    print("Простих чисел немає") 
else:
    # Обчислюємо та виводимо середнє арифметичне (сума/кількість)
    print("Середнє арифметичне простих:", S / k)
\end{py}
Програма почергово приймає цілі числа, поки не буде введено 0 (сигнал зупинки). Для кожного введеного числа перевіряється, чи є воно простим (тобто більшим за 1 і ділиться лише на 1 та саму себе).  Якщо число виявилося простим, програма збільшує лічильник $k$ на 1 та додає це число до загальної суми $S$. Після введення 0 програма ділить суму простих чисел на їхню кількість ($S$ / $k$) або повідомляє, що простих чисел не було введено. 
\begin{py}{Приклад 8. Перевірка наявності цифри 4 у натуральному числі}
N = int(input('N = '))  # Зчитуємо натуральне число
temp = N                # Створюємо тимчасову копію числа,
# щоб не змінювати оригінальне N

# Цикл триває, поки у тимчасовій змінній залишаються цифри 
#(число більше 0)
while temp > 0:
    x = temp % 10   # Отримуємо останню цифру числа 
                    #(остача від ділення на 10)

    if x == 4:      # Перевіряємо, чи дорівнює ця цифра 4
        print('Є 4')#Виводимо повідомлення, що цифру 4 знайдено
        break   # Достроково перериваємо цикл, 
                # оскільки перевіряти далі немає потреби
    #Відкидаємо останню цифру 
    temp //= 10) #(цілочисельне ділення на 10)

else: # Блок else виконується лише тоді,
   # коли цикл while завершився природним шляхом
   # (тобто temp стало <= 0), і не був перерваний break
   print('Немає 4')
\end{py}
Користувач вводить число N. Програма створює його копію temp, щоб зберегти оригінальне значення. Отримує останню цифру числа за допомогою temp \% 10.   Якщо цифра дорівнює 4, виводиться 'Є 4', і цикл негайно зупиняється (break).   Якщо цифра не 4, вона відкидається за допомогою цілочисельного ділення temp на 10. Якщо цикл закінчився і цифру 4 так і не знайшли, спрацьовує блок else, який виводить 'Немає 4'.   
\begin{py}{Приклад 9. Алгоритм Евкліда пошуку найбільшого спільного дільника (НСД)}
a = int(input("a = "))
b = int(input("b = "))

while a > 0 and b > 0:
    if a > b:
        a %= b
    else:
        b %= a

print("НСД =", a + b)
\end{py}
	\section*{Контрольні питання  }

\begin{enumerate}
	\item Який синтаксис циклу \texttt{while} у Python?
	\item У чому полягає відмінність між циклами \texttt{for} та \texttt{while}?
	\item Коли доцільно використовувати цикл \texttt{while}, а коли цикл \texttt{for}?
	\item Що означає «цикл із передумовою»?
	\item Що відбувається, якщо умова циклу \texttt{while} від самого початку має значення \texttt{False}?
	\item Як працюють оператори \texttt{break}, \texttt{continue}, \texttt{pass} у циклі \texttt{while}?
	\item Для чого використовується блок \texttt{else} у циклі \texttt{while}?
	\item Що таке «зациклювання»? Які основні причини його виникнення?
	\item Як уникнути нескінченного циклу у програмі?
	\item Наведіть приклад задачі, де цикл \texttt{while} зручніший за цикл \texttt{for}.
\end{enumerate}



\section*{Практичні завдання}

\begin{enumerate}
	\item Надрукувати числа від 0 до 9 за допомогою циклу \texttt{while}.
	\item Обчислити суму чисел від 1 до $n$, використовуючи цикл \texttt{while}.
	\item Запитувати у користувача пароль доти, доки він не введе правильний.
	\item Вивести таблицю квадратів чисел від 1 до 10.
	\item Реалізувати програму, яка знаходить найбільший спільний дільник двох чисел (алгоритм Евкліда).
	\item Написати програму, яка виводить усі цифри введеного числа по одній.
	\item Обчислити факторіал числа $n$ за допомогою циклу \texttt{while}.
	\item Запитувати у користувача додатне число доти, доки він не введе коректне значення.
	\item Обчислити середнє арифметичне простих чисел, що вводяться користувачем, до введення нуля.
	\item Реалізувати програму, яка перевіряє, чи містить введене число цифру нуль.
	\item Обчислити суму введених чисел до введення нуля (стоп‑сигнал).
	\item Написати нескінченний цикл \texttt{while True}, який виводить повідомлення, і перервати його комбінацією клавіш \texttt{Ctrl+C}.
\end{enumerate}


	\begin{thebibliography}{99}
	
	\bibitem{vasyliev}
	Васильєв О. \textit{Програмування мовою Python}. – Тернопіль: НК «Богдан», 2019. – 504 с.
	
	\bibitem{rudenko}
	Руденко В., Жугастров О. \textit{Інформатика. Основи алгоритмізації та програмування мовою Python}. – Харків: Ранок, 2019. – 192 с.
	
	\bibitem{kormen}
	Кормен Т., Лейзерсон Ч., Ривест Р. \textit{Вступ до алгоритмів} / Т. Кормен, Ч. Лейзерсон, Р. Ривест. – Київ: К.І.С., 2019. – 1288 с.
	
	\bibitem{matvienko}
	Матвієнко М. П. \textit{Теорія алгоритмів}. – Київ: Видавництво Ліра-К, 2019. – 344 с.
	
	\bibitem{borodkina}
	Бородкіна І. Л., Бородкін Г. О. \textit{Теорія алгоритмів. Посібник для студентів вищих навчальних закладів} / І. Л. Бородкіна, Г. О. Бородкін. – Київ: Центр навчальної літератури, 2019. – 184 с.
	
	\bibitem{matvienko-shapovalov}
	Матвієнко М. П., Шаповалов С. П. \textit{Математична логіка та теорія алгоритмів}. – Київ: Центр навчальної літератури, 2015. – 212 с.
	
\end{thebibliography}

	
	% === СПИСОК ЛІТЕРАТУРИ ===
	% Додайте бібліографічні посилання тут, використовуючи, наприклад, пакет biblatex/natbib
	% \bibliographystyle{plainnat}
	% \bibliography{references} 
	
	
\end{document}